% =============================================================================
%  lec17.tex —— 第 17 讲 导出 ∞-范畴（底本 PDF 第 131–138 页）
% =============================================================================
\lecture{17}{导出 $\infty$-范畴}
\begin{center}\small\color{black!55}（底本 PDF 第 131–138 页）\end{center}

本讲介绍两个稳定 $\infty$-范畴的例子：
\begin{itemize}[leftmargin=2.2em,itemsep=0.1em]
\item 加性范畴 $A$ 中上链复形的朴素 $\infty$-范畴 $\mathrm{K}(A)$；
\item Abel 范畴 $A$ 的导出 $\infty$-范畴 $\mathrm{D}(A)$。
\end{itemize}

\section{通过局部化构造 $\mathrm{K}(A)$ 与 $\mathrm{D}(A)$}

\begin{definition}\label{L17-1.1}
设 $A$ 是加性范畴，$\Ch(A)$ 是 $A$ 中上链复形的通常范畴。设 $f,g\colon M^{\bullet}\to N^{\bullet}$ 是 $\Ch(A)$ 中的态射。从 $f$ 到 $g$ 的一个\textbf{上链同伦}是一族映射 $h^n\colon M^n\to N^{n-1}$，使得对任意 $n$，
\[
f^n-g^n=d\circ h^n+h^{n+1}\circ d.
\]
若 $f$ 与 $g$ 之间存在上链同伦，则称 $f$ 与 $g$ \textbf{同伦}，或记作 $f\sim g$。
\end{definition}

\begin{remark}\label{L17-1.2}
同伦关系是 $\Hom_{\Ch(A)}(M^{\bullet},N^{\bullet})$ 上的等价关系。
\end{remark}

\begin{example}\label{L17-1.3}
设 $F,G\colon X\to Y$ 是拓扑空间之间的连续映射。则从 $F$ 到 $G$ 的任一同伦诱导从 $F^*$ 到 $G^*$ 的上链同伦，其中 $F^*\colon C^{\bullet}(Y)\to C^{\bullet}(X)$ 是奇异上链复形之间的上链同态。
\end{example}

\begin{exercise}\label{L17-1.4}
设 $A$ 是 Abel 范畴（从而可以定义上同调）。若 $f\sim g$，则作为态射 $H^n(M^{\bullet})\to H^n(N^{\bullet})$，有 $H^n(f)=H^n(g)$。
\end{exercise}

\begin{definition}\label{L17-1.5}
设 $A$ 是加性范畴，$f\colon M^{\bullet}\to N^{\bullet}$ 是 $\Ch(A)$ 中的态射。若存在 $g\colon N^{\bullet}\to M^{\bullet}$ 使得
\[
f\circ g\sim\id_{N^{\bullet}}\quad\text{且}\quad g\circ f\simeq\id_{M^{\bullet}},
\]
则称 $f$ 是\textbf{上链同伦等价}。设 $S$ 是 $\Ch(A)$ 中上链同伦等价构成的类。
\end{definition}

\begin{definition}\label{L17-1.6}
设 $A$ 是 Abel 范畴，$f\colon M^{\bullet}\to N^{\bullet}$ 是 $\Ch(A)$ 中的态射。若 $f$ 诱导上同调之间的同构，则称 $f$ 是\textbf{拟同构}。设 $W$ 是 $\Ch(A)$ 中拟同构构成的类。
\end{definition}

\begin{exercise}\label{L17-1.7}
设 $A$ 是 Abel 范畴。证明 $\Ch(A)$ 中的上链同伦等价是拟同构。
\end{exercise}

\begin{exercise}\label{L17-1.8}
设 $A$ 是 Abel 范畴，$0\to L\to M\to N\to 0$ 是短正合序列。证明
\[
\begin{tikzcd}[row sep=1.6em,column sep=1.7em]
\lbrack\cdots \arrow[r] & L \arrow[r] \arrow[d] & M \arrow[r] \arrow[d] & \cdots\rbrack \\
\lbrack\cdots \arrow[r] & 0 \arrow[r] & N \arrow[r] & \cdots\rbrack
\end{tikzcd}
\]
总是拟同构，但它是上链同伦等价当且仅当所给的短正合序列有分裂。
\end{exercise}

\begin{definition}\label{L17-1.9}
设 $A$ 是加性范畴。$A$ 中上链复形的\textbf{朴素 $\infty$-范畴}定义为 $\infty$-范畴局部化\footnote{按我们的约定，我们用同一符号表示通常范畴 $C$ 与相应的 $\infty$-范畴，后者实现为拟范畴 $\mathrm{N}^{\bullet}(C)$。因此 $\mathrm{K}(A)$ 可以表示为拟范畴局部化 $\mathrm{N}^{\bullet}(\Ch(A))[S^{-1}]$，这也是 Lurie 的书中使用的记号。}
\[
\mathrm{K}(A)\defeq\Ch(A)[S^{-1}].
\]
\end{definition}

\yizhu{“朴素”（naive）在此指未作局部化、直接以复形为对象并以上链同伦为态射的 $\infty$-范畴，与随后定义的“导出 $\infty$-范畴”相对；该词在中文文献中或译“天真”“朴素”，译本取“朴素”。\textbf{[来源：B]}}

\begin{definition}\label{L17-1.10}
设 $A$ 是 Abel 范畴。$A$ 的\textbf{无界导出 $\infty$-范畴}定义为 $\infty$-范畴局部化
\[
\mathrm{D}(A)\defeq\Ch(A)[W^{-1}].
\]
\end{definition}

\begin{exercise}\label{L17-1.11}
把 $\mathrm{D}(A)$ 写成 $\mathrm{K}(A)$ 的一个局部化。
\end{exercise}

\begin{warning}\label{L17-1.12}
在非 $\infty$-范畴的文献中，$\mathrm{K}(A)$ 与 $\mathrm{D}(A)$ 通常指 $1$-范畴局部化，它们是我们这些 $\infty$-范畴的同伦范畴。
\end{warning}

\nitem\label{L17-1.13} 众所周知（见 [Sta24, Tag 05QI]），$\mathrm{hK}(A)$ 具有三角范畴的结构，使得对 $\Ch(A)$ 中任意 $f\colon M^{\bullet}\to N^{\bullet}$，在 $\mathrm{hK}(A)$ 中相应的对象之间有特异三角
\[
M^{\bullet}\xrightarrow{\ f\ }N^{\bullet}\to\operatorname{cone}(f)\to M^{\bullet}[1].
\]
此外，当 $A$ 是 Abel 范畴时，该三角结构与拟同构类相容，从而在 $1$-范畴局部化 $\mathrm{hD}(A)$ 上诱导三角结构。

\begin{warning}\label{L17-1.14}
范畴 $\Ch(A)$ 几乎绝不是三角范畴，因为任何 Abel 三角范畴都是半单的。
\end{warning}

\nitem\label{L17-1.15} 下述结果表明上述三角范畴来自稳定 $\infty$-范畴。

\begin{theorem}\label{L17-1.16}
（HA.1.3.2.10）设 $A$ 是加性范畴。$\infty$-范畴 $\mathrm{K}(A)$ 是稳定的，且 $\mathrm{hK}(A)$ 上的三角结构与熟知的三角结构一致\footnote{这一论断可由 HA.1.3.2.10 的证明推出。}。
\end{theorem}

\begin{theorem}\label{L17-1.17}
（HA.1.3.5.9, 1.3.5.21）设 $A$ 是 Grothendieck\footnote{这是指 $A$ 可表现，且取滤过余极限是正合的。} Abel 范畴。$\infty$-范畴 $\mathrm{D}(A)$ 可表现且稳定，且 $\mathrm{hD}(A)$ 上的三角结构与熟知的三角结构一致。
\end{theorem}

\begin{warning}\label{L17-1.18}
设 $A$ 是 Grothendieck Abel 范畴。$\infty$-范畴 $\mathrm{K}(A)$ 几乎绝不是可表现的。见 MathOverflow 上的这个问题。
\end{warning}

\begin{remark}\label{L17-1.19}
注意，上面定义的 $\infty$-范畴 $\mathrm{D}(A)$ 对任意 Abel 范畴 $A$ 都有意义。可以证明 $\mathrm{D}(A)$ 总是稳定的（例如可用 [NS18, Theorem I.3.3]）。然而 $\mathrm{D}(A)$ 一般可能行为不良。例如，当 $A$ 小时，最好考虑 $\mathrm{D}^{\mathrm{b}}(A)$。
\end{remark}

\section{通过微分分次脉构造 $\mathrm{K}(A)$ 与 $\mathrm{D}(A)$}

\nitem\label{L17-2.1} 定义~\ref{L17-1.9} 与定义~\ref{L17-1.10} 简洁，但不便于计算。本节给出这些 $\infty$-范畴的一个更显式的构造。

\nitem\label{L17-2.2} $\mathrm{K}(A)$ 的另一种构造将是一个拟范畴，使得
\begin{itemize}[leftmargin=2.2em,itemsep=0.1em]
\item 对象是上链复形 $M_0$；
\item 态射是上链同态 $M_0\to M_1$；
\item 一个 $2$-单形由下图给出
\[
\begin{tikzcd}[row sep=2.2em,column sep=2.6em]
& M_1 \arrow[dr, "f_{12}"] \arrow[d, Rightarrow, "f_{012}"] & \\
M_0 \arrow[ur, "f_{01}"] \arrow[rr, "f_{02}"'] & {} & M_2
\end{tikzcd}
\]
其中 $h_{012}$ 表示从 $f_{12}\circ f_{01}$ 到 $f_{02}$ 的同伦；
\item $\cdots$
\end{itemize}
为描述高维单形，引入一些术语会比较方便。

\begin{construction}\label{L17-2.3}
设 $\Ch(\Ab)$ 是上链复形的通常范畴。对 $M,N\in\Ch(\Ab)$，\textbf{分次张量积} $M\otimes N$ 定义如下。
\begin{itemize}[leftmargin=2.2em,itemsep=0.1em]
\item 对每个整数 $k$，
\[
(M\otimes N)^k\defeq\bigoplus_{i+j=k}M^i\otimes N^j.
\]
\item 对每个整数 $k$，微分 $(M\otimes N)^k\to(M\otimes N)^{k+1}$ 由\textbf{分次 Leibniz 法则}给出，它把纯张量 $m\otimes n\in M^i\otimes N^j$ 送到
\[
d(m\otimes n)\defeq d(m)\otimes n+(-1)^j m\otimes d(n).
\]
\end{itemize}
$\Ch(\Ab)$ 上有自然的幺半结构，其乘法由分次张量积给出。
\end{construction}

\begin{example}\label{L17-2.4}
设 $X$、$Y$ 是拓扑空间。则 $C^{\bullet}(X\times Y)\simeq C^{\bullet}(X)\otimes C^{\bullet}(Y)$。
\end{example}

\begin{exercise}\label{L17-2.5}
设 $M\in\Ch(\Ab)$。证明 $M$ 中的 $0$-余圈等同于一个态射 $\mathbf{1}\to M$，其中 $\mathbf{1}$ 是幺半单位。
\end{exercise}

\begin{definition}\label{L17-2.6}
\textbf{微分分次范畴}，或 \textbf{dg-范畴}，是 $\Ch(\Ab)$-充实范畴。
\end{definition}

\nitem\label{L17-2.7} 按定义，一个 dg-范畴 $\mathbb{C}$ 由以下数据组成：
\begin{itemize}[leftmargin=2.2em,itemsep=0.1em]
\item 对对象 $x,y\in\mathbb{C}$，一个上链复形 $\Hom_{\mathbb{C}}(x,y)\in\Ch(\Ab)$，称为\textbf{映射复形}；
\item 一个由上链同态给出的复合法则
\[
-\circ-\colon\Hom_{\mathbb{C}}(x,y)\otimes\Hom_{\mathbb{C}}(y,z)\to\Hom_{\mathbb{C}}(x,z),
\]
它在显然的意义下结合；
\item 一个 $0$-余圈 $\id_x\in\Hom_{\mathbb{C}}(x,x)^0$，使得对任意 $f\in\Hom_{\mathbb{C}}(x,y)^m$ 与 $g\in\Hom_{\mathbb{C}}(y,x)^n$，有 $f\circ\id_x=f$ 与 $\id_x\circ g=g$。
\end{itemize}

\begin{exercise}\label{L17-2.8}
证明 $\Hom_{\mathbb{C}}(x,x)^0$ 中满足上述条件的任意 $\id_x$ 自动是余圈。
\end{exercise}

\begin{definition}\label{L17-2.9}
设 $\mathbb{C}$ 是 dg-范畴。$\mathbb{C}$ 的\textbf{底层范畴} $\mathrm{C}$ 定义为这样的范畴：其 Hom 集由
\[
\Hom_{\mathrm{C}}(x,y)\defeq\Hom_{\mathbb{C}\mathrm{h}(A)}(\mathbf{1},\Hom_{\mathbb{C}}(x,y))
\]
给出。
\end{definition}

\nitem\label{L17-2.10} 我们将把 $\Hom_{\mathrm{C}}(x,y)$ 与 $\Hom_{\mathbb{C}}(x,y)$ 中 $0$-余圈构成的 Abel 群等同。

\begin{definition}\label{L17-2.11}
设 $\mathbb{C}$ 是 dg-范畴。$\mathbb{C}$ 的\textbf{同伦范畴} $\mathrm{h}\mathbb{C}$ 定义为这样的范畴：其 Hom 集由
\[
\Hom_{\mathrm{h}\mathbb{C}}(x,y)\defeq H^0(\Hom_{\mathbb{C}}(x,y))
\]
给出。
\end{definition}

\begin{example}\label{L17-2.12}
对任意加性范畴 $A$，通常范畴 $\Ch(A)$ 有到 $\Ch(\Ab)$ 上的自然充实 $\mathbb{C}\mathrm{h}(A)$，使得
\[
\Hom_{\mathbb{C}\mathrm{h}(A)}(M,N)^k\defeq\prod_i\Hom_A(M^i,N^{i+k}),
\]
其微分由分次 Leibniz 法则
\[
(df)(x)=d(f(x))-(-1)^k f(dx)
\]
给出，其中 $f\in\Hom_{\mathbb{C}\mathrm{h}(A)}(M,N)^k$，$x\in M^i$。
\end{example}

\begin{exercise}\label{L17-2.13}
证明 $\mathbb{C}\mathrm{h}(A)$ 的底层范畴确实是 $\Ch(A)$。它的同伦范畴是什么？
\end{exercise}

\begin{exercise}\label{L17-2.14}
设 $f,g\in\Hom_{\mathbb{C}\mathrm{h}(A)}(M,N)^0$ 是 $0$-余圈，视为通常范畴 $\Ch(A)$ 中的态射。证明从 $f$ 到 $g$ 的同伦等同于一个元素 $h\in\Hom_{\mathbb{C}\mathrm{h}(A)}(M,N)^{-1}$ 使得 $dh=f-g$。
\end{exercise}

\nitem\label{L17-2.15}\textbf{命题—构造}（HA.1.3.1.10）\quad 设 $\mathbb{C}$ 是 dg-范畴。我们有一个拟范畴 $\mathrm{N}_{dg}(\mathbb{C})$，称为 $\mathbb{C}$ 的 \textbf{dg-脉}，定义如下。
\begin{itemize}[leftmargin=2.2em,itemsep=0.1em]
\item 一个 $0$-单形是 $\mathbb{C}$ 中的对象；
\item 一个 $1$-单形由其边界 $\{X_0,X_1\}$ 连同一个 $0$-余圈 $f_{01}\in\Hom_{\mathbb{C}}(X_0,X_1)^0$ 组成；
\item 一个 $2$-单形由其边界 $\{f_{01},f_{02},f_{12}\}$ 连同一个元素 $g_{012}\in\Hom_{\mathbb{C}}(X_0,X_2)^{-1}$ 组成，使得
\[
df_{012}=-(f_{02}-f_{12}\circ f_{01}).
\]
\item 一个 $3$-单形由其边界 $\{f_{012},f_{013},f_{023},f_{123}\}$ 连同一个元素 $f_{0123}\in\Hom_{\mathbb{C}}(X_0,X_3)^{-2}$ 组成，使得
\[
df_{0123}=-(f_{013}-f_{23}\circ f_{012})+(f_{023}-f_{123}\circ f_{01}).
\]
\item 更高维的单形类似地给出（见 HA.1.3.1.6）。
\end{itemize}

\yizhu{原文的记号有几处前后不一致，译文均照录：2.2 段的图中 $2$-单形的填充子标作 $f_{012}$，正文却称它为 $h_{012}$；本构造中 $2$-单形所附的元素记作 $g_{012}$，条件却写作 $df_{012}=\cdots$；此外，定义~\ref{L17-2.9} 的 $\Hom_{\mathbb{C}\mathrm{h}(A)}$ 中出现的 $A$ 在该处并未引入（按上下文应为 $\Ch(\Ab)$）。\textbf{[来源：I]}}

\begin{proposition}\label{L17-2.16}
设 $\mathbb{C}$ 是 dg-范畴，$X,Y$ 是 $\mathbb{C}$ 中的对象。则有典范同构
\[
\pi_i(\mathrm{Maps}_{\mathrm{N}_{dg}(\mathbb{C})}(X,Y))\simeq H^{-i}(\Hom_{\mathbb{C}}(X,Y)).
\]
\end{proposition}

\begin{remark}\label{L17-2.17}
我们下次证明上述命题。事实上，空间 $\mathrm{Maps}_{\mathrm{N}_{dg}(\mathbb{C})}(X,Y)$ 可以由截断 $\tau^{\le0}\Hom_{\mathbb{C}}(X,Y)$ 以典范方式得到，这称为 Dold–Kan 对应。
\end{remark}

\begin{exercise}\label{L17-2.18}
构造一个单态射 $\mathrm{N}(\mathrm{C})\to\mathrm{N}_{dg}(\mathbb{C})$，它在 $n\le1$ 的 $n$-单形上是双射，其中 $\mathrm{N}(\mathrm{C})$ 是 $\mathbb{C}$ 的底层通常范畴 $\mathrm{C}$ 的通常脉。
\end{exercise}

\begin{exercise}\label{L17-2.19}
构造一个等价 $\mathrm{hN}_{dg}(\mathbb{C})\simeq\mathrm{h}\mathbb{C}$。
\end{exercise}

\begin{proposition}\label{L17-2.20}
（HA.1.3.4.5）设 $A$ 是加性范畴。则函子
\[
\mathrm{N}(\Ch(A))\to\mathrm{N}_{dg}(\mathbb{C}\mathrm{h}(A))
\]
把 $\mathrm{N}_{dg}(\mathbb{C}\mathrm{h}(A))$ 实现为 $\mathrm{N}(\Ch(A))$ 关于上链同伦等价类的局部化。换言之，$\mathrm{N}_{dg}(\mathbb{C}\mathrm{h}(A))$ 表示 $\infty$-范畴 $\mathrm{K}(A)$。
\end{proposition}

\begin{exercise}\label{L17-2.21}
设 $A$ 是 Abel 范畴，$0\to L\to M\to N\to 0$ 是短正合序列。证明它表示 $\mathrm{K}(A)$ 中的零对象，当且仅当它允许一个分裂。
\end{exercise}

\nitem\label{L17-2.22} 我们的下一个目标，是在 $A$ 为 Grothendieck Abel 范畴时，把 $\mathrm{D}(A)$ 实现为某个 dg-范畴的 dg-脉。回忆这样的 $A$ 有足够多的内射对象。

\begin{theorem}\label{L17-2.23}
设 $A$ 是 Grothendieck Abel 范畴。则 $\Ch(A)$ 容许一个左真组合模型结构，由以下各项确定。
\begin{itemize}[leftmargin=3.4em,itemsep=0.1em]
\item[(W)] 弱等价是拟同构；
\item[(C)] 上纤维化是逐次单态射。
\end{itemize}
\end{theorem}

\begin{exercise}\label{L17-2.24}
若 $M\in\Ch(A)$ 在上述模型结构中是纤维对象，则每个 $M^n\in A$ 都是内射对象。反之，若 $M$ 下有界且每个 $M^n$ 都是内射对象，则 $M$ 是纤维对象。
\end{exercise}

\begin{proposition}\label{L17-2.25}
（HA.1.3.5.13）设 $\Ch(A)^{\circ}\subseteq\Ch(A)$ 是双纤维对象构成的全子范畴。设 $\mathbb{C}\mathrm{h}(A)^{\circ}\subseteq\mathbb{C}\mathrm{h}(A)$ 是相应的 dg-范畴。则嵌入
\[
\mathrm{N}_{dg}(\mathbb{C}\mathrm{h}(A)^{\circ})\to\mathrm{N}_{dg}(\mathbb{C}\mathrm{h}(A))
\]
有左伴随
\[
\mathrm{N}_{dg}(\mathbb{C}\mathrm{h}(A))\to\mathrm{N}_{dg}(\mathbb{C}\mathrm{h}(A)^{\circ}),
\]
使得复合
\[
\mathrm{N}(\Ch(A))\to\mathrm{N}_{dg}(\mathbb{C}\mathrm{h}(A))\to\mathrm{N}_{dg}(\mathbb{C}\mathrm{h}(A)^{\circ})
\]
把 $\mathrm{N}_{dg}(\mathbb{C}\mathrm{h}(A)^{\circ})$ 实现为 $\mathrm{N}(\Ch(A))$ 关于拟同构类的局部化。换言之，$\mathrm{N}_{dg}(\mathbb{C}\mathrm{h}(A)^{\circ})$ 表示 $\infty$-范畴 $\mathrm{D}(A)$。
\end{proposition}

\nitem\label{L17-2.26} 上述命题蕴含：对任意 Grothendieck Abel 范畴 $A$，有正合函子之间的伴随
\begin{equation}\label{L17-eq2.1}
\begin{tikzcd}[column sep=3.4em]
\mathrm{K}(A) \arrow[r, shift left=0.7ex] & \mathrm{D}(A) \arrow[l, shift left=0.7ex]
\end{tikzcd}
\end{equation}
使得右伴随函子全忠实，且左伴随由取纤维替换给出。

\begin{exercise}\label{L17-2.27}
求左伴随函子 $\mathrm{K}(A)\to\mathrm{D}(A)$ 的核。
\end{exercise}

\section{t-结构}

\begin{proposition}\label{L17-3.1}
（HA.1.3.5.19, 1.3.5.21）设 $A$ 是 Grothendieck Abel 范畴。
\begin{itemize}[leftmargin=2.2em,itemsep=0.1em]
\item 设 $\mathrm{D}(A)^{\le0}$ 是由满足 $H^n(M)\simeq0$（$n>0$）的上链复形 $M$ 所表示的对象构成的全子 $\infty$-范畴。
\item 设 $\mathrm{D}(A)^{\ge0}$ 是由满足 $H^n(M)\simeq0$（$n<0$）的上链复形 $M$ 所表示的对象构成的全子 $\infty$-范畴。
\end{itemize}
则 $(\mathrm{D}(A)^{\le0},\mathrm{D}(A)^{\ge0})$ 确定 $\mathrm{D}(A)$ 上的一个 t-结构。此外，
\begin{itemize}[leftmargin=2.2em,itemsep=0.1em]
\item 函子 $H^0\colon\Ch(A)\to A$ 诱导等价 $\mathrm{D}(A)^{\heartsuit}\to A$。
\item 该 t-结构可及、左分离、右完备，且与滤过余极限相容。
\end{itemize}
\end{proposition}

\begin{warning}\label{L17-3.2}
$\mathrm{D}(A)$ 上的 t-结构一般不是左完备的。反例见 [Nee11]。左完备化常记作 $\widehat{\mathrm{D}}(A)$。
\end{warning}

\begin{example}\label{L17-3.3}
设 $A=\Mod_R^{\heartsuit}$ 是结合环 $R$ 的 $R$-模范畴。在以后的讲中，我们将构造
\begin{itemize}[leftmargin=2.2em,itemsep=0.1em]
\item $\mathrm{Sptr}$ 上由 smash 积给出的对称幺半结构；
\item Eilenberg–MacLane 谱 $\mathbb{H}R\in\mathrm{Sptr}$ 上的 $E_1$-代数结构；
\item $\mathrm{Sptr}$ 中 $\mathbb{H}R$-模的稳定 $\infty$-范畴 $\Mod_{\mathbb{H}R}$，配有从 $\mathrm{Sptr}$ 的 t-结构诱导的 t-结构；
\item 一个 t-正合等价
\[
\Mod_{\mathbb{H}R}\simeq\mathrm{D}(\Mod_R^{\heartsuit})\eqqcolon\mathrm{D}(R).
\]
\end{itemize}
此外，当 $R$ 交换时，我们将给 $\mathrm{D}(R)$ 配备由 $\mathbb{H}R$ 上的相对张量积给出的对称幺半结构。
\end{example}

\begin{exercise}\label{L17-3.4}
设 $A$ 是 Grothendieck Abel 范畴。证明 $\mathrm{D}(A)$ 的左有界部分可以等同于
\[
\mathrm{D}^+(A)\simeq\mathrm{N}_{dg}(\mathbb{C}\mathrm{h}^+(A_{\mathrm{inj}})),
\]
其中 $\mathbb{C}\mathrm{h}^+(A_{\mathrm{inj}})$ 是 $A$ 中内射对象的左有界复形构成的 dg-范畴。注意，该证明对 $\mathrm{D}^-(A)$ 应当不成立。
\end{exercise}

\nitem\label{L17-3.5} 受上述练习的启发，我们作如下定义。

\begin{definition}\label{L17-3.6}
设 $A$ 是有足够多内射对象的 Abel 范畴。$A$ 的\textbf{左有界导出 $\infty$-范畴}定义为
\[
\mathrm{D}^+(A)\defeq\mathrm{N}_{dg}(\mathbb{C}\mathrm{h}^+(A_{\mathrm{inj}})).
\]
对偶地，设 $A$ 是有足够多投射对象的 Abel 范畴。$A$ 的\textbf{右有界导出 $\infty$-范畴}定义为
\[
\mathrm{D}^-(A)\defeq\mathrm{N}_{dg}(\mathbb{C}\mathrm{h}^-(A_{\mathrm{proj}}))\simeq\mathrm{D}^+(A^{\op})^{\op}.
\]
\end{definition}

\yizhu{作者把 $\mathrm{D}^+(A)$（复形下有界）称为“左有界”、把 $\mathrm{D}^-(A)$（复形上有界）称为“右有界”，这与常见文献的“下有界/上有界”说法不同，但所指相同；译文照原文的方向词。\textbf{[来源：B]}}

\begin{exercise}\label{L17-3.7}
证明 $\mathrm{D}^-(A)$ 是 $\mathrm{K}(A_{\mathrm{proj}})$ 的稳定子 $\infty$-范畴。
\end{exercise}

\begin{warning}\label{L17-3.8}
$\infty$-范畴 $\mathrm{D}^+(A)$ 与 $\mathrm{D}^-(A)$ 不可表现，因为可数余积可能不存在。
\end{warning}

\nitem\label{L17-3.9} 下述结果用泛性质刻画 $\mathrm{D}^-(A)$。

\begin{proposition}\label{L17-3.10}
（HA.1.3.2.19, 1.3.3.16）设 $A$ 是有足够多投射对象的 Abel 范畴。显然的选择在 $\mathrm{D}^-(A)$ 上定义了一个左完备 t-结构，其心典范等价于 $A$。
\end{proposition}

\begin{theorem}\label{L17-3.11}
（HA.1.3.3.2, 1.3.3.6）设 $A$ 是有足够多投射对象的 Abel 范畴。对任意配有左完备 t-结构的稳定 $\infty$-范畴 $\cC$，以下数据等价：
\begin{enumerate}[label=(\roman*),leftmargin=2.4em,itemsep=0.1em]
\item 一个右 t-正合函子 $F\colon\mathrm{D}^-(A)\to\cC$，它把 $A_{\mathrm{proj}}$ 送入 $\cC^{\heartsuit}$；
\item 一个右正合函子 $f\colon A\to\cC^{\heartsuit}$，
\end{enumerate}
其中 $f\defeq\tau^{\ge0}\circ F|_{A}$。此外，以下条件等价
\begin{itemize}[leftmargin=2.2em,itemsep=0.1em]
\item $F$ 是 t-正合的；
\item $F$ 保持心；
\item $f$ 是正合的。
\end{itemize}
\end{theorem}

\begin{definition}\label{L17-3.12}
设 $A$ 是有足够多投射对象的 Abel 范畴，$\cC$ 是配有左完备 t-结构的稳定 $\infty$-范畴。对右正合函子 $f\colon A\to\cC^{\heartsuit}$，相应的函子 $F$ 称为 $f$ 的\textbf{左导出函子}，并记作
\[
\mathbb{L}f\colon\mathrm{D}^-(A)\to\cC.
\]
\end{definition}

\begin{exercise}\label{L17-3.13}
证明 $\mathbb{L}f(M)$ 可以通过复形 $M$ 的投射消解来计算。
\end{exercise}

\begin{exercise}\label{L17-3.14}
考虑对应于等价 $\Ab\simeq\mathrm{Sptr}^{\heartsuit}$ 的 t-正合函子 $\mathrm{D}^-(\Ab)\to\mathrm{Sptr}$，$A\mapsto HA$。取右完备化，得到函子
\[
\pi_*\colon\mathrm{D}(\Ab)\to\mathrm{Sptr}.
\]
证明：
\begin{enumerate}[label=(\arabic*),leftmargin=2.4em,itemsep=0.1em]
\item 该函子也可以通过定理~\ref{L17-3.11} 的对偶所描述的 $\mathrm{D}^+(\Ab)$ 的泛性质得到。
\item 猜测复合 $\Omega^{\infty}\circ\pi_*$。
\end{enumerate}
\end{exercise}

\nitem\label{L17-3.15} 设 $A$ 是有足够多投射对象的 Grothendieck Abel 范畴。下述结果表明 $\mathrm{D}^-(A)$ 的两种定义一致。

\begin{proposition}\label{L17-3.16}
（HA.1.3.5.24）设 $A$ 是有足够多投射对象的 Grothendieck Abel 范畴。则
\[
\mathrm{D}^-(A)\xrightarrow{\ \ \subseteq\ \ }\mathrm{K}(A)\to\mathrm{D}(A)
\]
全忠实，且本质像由 $\mathrm{D}(A)$ 的右有界部分给出。特别地，$\mathrm{D}(A)$ 左完备。
\end{proposition}

\nitem\label{L17-3.17} 事实上，当 $A$ 是 Grothendieck Abel 范畴时，$\mathrm{K}(A)$ 上也有典范的 t-结构。然而，需要小心余连通部分，因为 $\mathrm{K}(A)$ 中存在上同调为零的非零对象。

\begin{proposition}\label{L17-3.18}
（HA.1.3.5.18）设 $A$ 是 Grothendieck Abel 范畴。
\begin{itemize}[leftmargin=2.2em,itemsep=0.1em]
\item 设 $\mathrm{K}(A)^{\le0}$ 是由满足 $H^n(M)\simeq0$（$n>0$）的上链复形 $M$ 所表示的对象构成的全子 $\infty$-范畴。
\item 设 $\mathrm{K}(A)^{\ge0}$ 是由满足 $M^n\simeq0$（$n<0$）且每个 $M^n$ 都是内射对象的上链复形 $M$ 所表示的对象构成的全子 $\infty$-范畴。
\end{itemize}
则 $(\mathrm{K}(A)^{\le0},\mathrm{K}(A)^{\ge0})$ 确定 $\mathrm{K}(A)$ 上的一个 t-结构。
\end{proposition}

\begin{exercise}\label{L17-3.19}
把 $\mathrm{K}(A)^{\ge0}$ 等同于全忠实函子
\[
\mathrm{D}(A)^{\ge0}\to\mathrm{D}(A)\to\mathrm{K}(A)
\]
的本质像，其中第一个函子是式~\eqref{L17-eq2.1} 中的右伴随。
\end{exercise}

\begin{exercise}\label{L17-3.20}
证明 $\mathrm{K}(A)\rightleftarrows\mathrm{D}(A)$ 中的两个函子都是 t-正合的，且它们诱导心之间的等价。
\end{exercise}

\begin{exercise}\label{L17-3.21}
证明 $\mathrm{K}(A)$ 上的 t-结构右完备且与滤过余极限相容，但不是左分离的。
\end{exercise}

\appendixsection{$\widetilde{\mathrm{D}}(A)$}

\begin{construction}\label{L17-A.1}
设 $A$ 是 Grothendieck Abel 范畴。定义 $A$ 的\textbf{非分离导出 $\infty$-范畴}为
\[
\widetilde{\mathrm{D}}(A)\defeq\mathrm{N}^{dg}(\mathbb{C}\mathrm{h}(A_{\mathrm{inj}}))\simeq\mathrm{K}(A_{\mathrm{inj}}).
\]
$\widetilde{\mathrm{D}}(A)$ 上有一个类似于 $\mathrm{K}(A)$ 上的 t-结构，使得
\begin{itemize}[leftmargin=2.2em,itemsep=0.1em]
\item 心典范地等同于 $A$；
\item 该 t-结构可及、右完备，且与滤过余极限相容；
\item 该 t-结构是（左）反完备的。
\end{itemize}
\end{construction}

\begin{remark}\label{L17-A.2}
大致说来，反完备意味着该 t-结构与完备的 t-结构“正交”。事实上，存在本质唯一的保持余极限的 t-正合函子
\[
\widetilde{\mathrm{D}}(A)\to\mathrm{D}(A),
\]
它限制到心上为恒等函子。此外，该函子把 $\widetilde{\mathrm{D}}(A)$ 实现为 $\mathrm{D}(A)$ 的（左）反完备化。
\end{remark}

\begin{example}\label{L17-A.3}
对 Noetherian 交换环 $R$，
\[
\widetilde{\mathrm{D}}(\Mod_R^{\heartsuit})\simeq\Ind(\mathrm{D}^{\mathrm{b}}(\Mod_{R,\mathrm{fg}}^{\heartsuit})),
\]
其中 $\mathrm{D}^{\mathrm{b}}(\Mod_{R,\mathrm{fg}}^{\heartsuit})$ 是 $\mathrm{D}(\Mod_R^{\heartsuit})$ 中由上同调有界且有限生成的上链复形构成的全子 $\infty$-范畴。
\end{example}

\nitem\label{L17-A.4} \textbf{建议阅读}\quad SAG.C。

\appendixsection{dg-范畴 vs. $\mathbb{H}Z$-线性稳定 $\infty$-范畴}

\nitem\label{L17-B.1} 任意 dg-范畴 $\mathbb{C}$ 的同伦范畴 $\mathrm{h}\mathbb{C}$ 有一个三角结构的典范候选。若该候选确实是三角结构，则称 $\mathbb{C}$ 是\textbf{预三角}的。以下概念本质等价：
\begin{itemize}[leftmargin=2.2em,itemsep=0.1em]
\item 交换环 $R$ 上的预三角 dg-范畴；
\item $\mathbb{H}R$-线性稳定 $\infty$-范畴。
\end{itemize}

\nitem\label{L17-B.2} \textbf{建议阅读}\quad [Coh13]。

\subsection*{参考文献}
\par\noindent\hangindent=5.2em\hangafter=1 [Coh13]\hspace{0.6em}Lee Cohn. \emph{Differential graded categories are k-linear stable infinity categories}. arXiv preprint arXiv:1308.2587, 2013.
\par\noindent\hangindent=5.2em\hangafter=1 [Nee11]\hspace{0.6em}Amnon Neeman. \emph{Non-left-complete derived categories}. \emph{Mathematical Research Letters}, 18(5):827--832, 2011.
\par\noindent\hangindent=5.2em\hangafter=1 [NS18]\hspace{0.6em}Thomas Nikolaus and Peter Scholze. \emph{On topological cyclic homology}. \emph{Acta Math}, 221:203--409, 2018.
\par\noindent\hangindent=5.2em\hangafter=1 [Sta24]\hspace{0.6em}The Stacks project authors. \emph{The stacks project}. \url{https://stacks.math.columbia.edu}, 2024.
